% Chapter 3, direct and inverse problems (Sections 3.2-3.3), in the style of
% fig:3Dbody. (a) complementary data: u = u^D on S_u, flux on S_phi, crack
% known. (b) overdetermined data: u and its normal derivative known on the
% whole boundary (two offset contours), crack unknown.
\begin{tikzpicture}[scale=0.9,>={Latex[length=2mm]}]
\def\blob{(A) to[out=-70,in=180] (B) to[out=0,in=-130] (C)
    to[out=50,in=-30] (D) to[out=150,in=10] (E) to[out=190,in=110] (A)}
\tikzset{body/.style={fill=gray!12},
  crack/.style={line width=1.8pt,line cap=round}}
% ---------------- (a) direct problem ----------------
\begin{scope}
  \coordinate (A) at (0,0.6);   \coordinate (B) at (2.0,-0.4);
  \coordinate (C) at (4.4,0.5); \coordinate (D) at (3.9,2.6);
  \coordinate (E) at (1.0,2.5);
  % S_u (electrodes, blue) and S_phi (flux, green)
  \draw[cu!30,line width=6pt] (E) to[out=190,in=110] (A) to[out=-70,in=180] (B);
  \fill[body] \blob -- cycle;
  \draw[cu,very thick] (E) to[out=190,in=110] (A) to[out=-70,in=180] (B);
  \draw[cphi,very thick,postaction={decorate},decoration={markings,
     mark=at position 0.80 with {\draw[->,thick,black] (0,0) -- (0,-0.75) coordinate (N);}}]
    (B) to[out=0,in=-130] (C) to[out=50,in=-30] (D) to[out=150,in=10] (E);
  \node[left] at (N) {$\vnu$};
  % the crack, known
  \draw[crack] (1.75,1.05) -- (2.85,1.6);
  \node[below right=-1pt] at (2.55,1.45) {$\Gamma$};
  \node at (1.35,1.75) {$\Delta u=0$};
  \node[cu,anchor=north east,align=right,font=\small] at (1.05,-0.4)
    {$S_u$: $u=\uD$};
  \node[cphi,anchor=west,align=left,font=\small] at (4.75,1.75)
    {$S_\varphi$:\\ $\partial_\nu u=\pD$};
  \node[font=\sffamily] at (2.2,-1.35) {(a) direct problem};
\end{scope}
% ---------------- (b) inverse problem ----------------
\begin{scope}[xshift=8.0cm]
  \coordinate (A) at (0,0.6);   \coordinate (B) at (2.0,-0.4);
  \coordinate (C) at (4.4,0.5); \coordinate (D) at (3.9,2.6);
  \coordinate (E) at (1.0,2.5);
  % two offset contours: u known (blue, inner), flux known (green, outer)
  \draw[cphi,line width=8pt,line join=round] \blob -- cycle;
  \draw[white,line width=5.4pt,line join=round] \blob -- cycle;
  \draw[cu,line width=3.2pt,line join=round] \blob -- cycle;
  \fill[body] \blob -- cycle;
  \draw[gray!50!black,thin] \blob -- cycle;
  % the crack, unknown
  \draw[crack,red!70!black,dash pattern=on 4pt off 2.5pt,line cap=butt]
    (1.75,1.05) -- (2.85,1.6);
  \node[red!70!black,below right=-1pt] at (2.55,1.45) {$\Gamma\,?$};
  \node at (1.35,1.75) {$\Delta u=0$};
  \draw[cu,thin] (0.45,2.25) -- (-0.15,2.75) node[above left=-2pt,font=\small]
    {$u=\uD$ on $\partial\Omega$};
  \draw[cphi,thin] (3.95,-0.05) -- (4.5,-0.45) node[below right=-2pt,font=\small]
    {$\partial_\nu u=\pD$ on $\partial\Omega$};
  \node[font=\sffamily] at (2.2,-1.35) {(b) inverse problem};
\end{scope}
\end{tikzpicture}
